\documentclass{article}
\usepackage{graphicx} % Required for inserting images
\usepackage{xcolor}

\title{\textbf{Proposal:} \\ Formal Verification of Collusion in Transparent AI Agents}
\author{Colomban Duclaux}
\date{\today}

\usepackage[backend=biber,style=alphabetic]{biblatex} % Ensure style matches your preference
\addbibresource{references.bib} % Link your .bib file


\newcommand{\dupoc}{\texttt{DUPOC}}
\newcommand{\cupod}{\texttt{CUPOD}}
\newcommand{\cimcic}{\texttt{CIMCIC}}
\newcommand{\dimcid}{\texttt{DIMCID}}

\begin{document}

\maketitle

\section{Introduction}
We want to be able to better understand how AI agents coordinate/collude together or against each other. How they manage to do this depends on the information they have about each other, going from fully open-source to black-box.

\textbf{Question:} "How well does controlling 
the transparency level serve as a lever for preventing collusion in oversight mechanisms?"

To respond to this question, 3 concrete scenarios: \begin{enumerate}
    \item \textbf{ Collusion via credible commitments under full visibility} $\to$ White box: \begin{itemize}
        \item Implement formal verifier agents (CUPOD/DUPOC-like strategies)
        \item Examine the minimal paths that lead to AI agents forming collusive agreements, and characterize the conditions
        \item Investigate system designs and agents behaviors that can prevent collusion, using a proof assistant.
    \end{itemize}
    \item Collusion via reverse-engineering the opponent's objectives (simulation-based collusion) $\to$ Gray box: \begin{itemize}
        \item model weight access, reasoning traces, output-only, etc...
        \item How much transparency is the minimum needed to catch deceptive behaviour? (use of verification model for this)
        \item Can we trust what AIs claim as their own alignement?
    \end{itemize}
    \item Collusion via deceptive alignment (strategically hiding information) $\to$ Black box: \begin{itemize}
        \item What strategies emerge in this scenario for collusion?
        \item Is AI reasoning powerful enough that collusion arises automatically?
    \end{itemize}
\end{enumerate}

\section{Methodology}
\begin{itemize}
    \item White box: Full access to one another's source code. Each agent is instantiated as a Lean-verifiable program to encode commitments, defections, and counterfactual reasoning as formal propositions.
    Introduction of mechanism-level constraints, including asymmetric visibility and probabilistic agents to evaluate the sweet spot between collusion and collaboration.
    
    Formulation of Collusion Dynamics Dataset using the Lean-verifiable agents to form upper bounds on collusion under maximal transparency.

    Use of Lean to formally prove upper bounds on achievable collusion rates. Empirical tournament results will be compared to these bounds.
    \item Gray box: Model the amount of information to disclose as a strategic choice pondered by the cost upon information disclosure. Compare results to phase 1.
    
    Formulation of the Strategic Revelation Dataset under varying transparency costs.

    Quantify the reduction in collusion frequency given a specific mechanism. Robustness tests using randomized perturbations and cost penalties for communication or simulations.
    \item Black box: LLM-based agents with recursive reasoning templates. Evaluate wether collusion against the overseer emerges implicitly as a function of modeling depth. Compare to phases 1 and 2, testing whether behavioral inference enables the same collusion patterns as in the other phases. Aims to see if the removal of transparency is sufficient to disrupt collusion.
    
    Formulation of the Reasoning Depths Dataset using LLM-based agents acting through behavioral inference, guided by fixed reasoning templates.

    Collusion rates analyzed as a function of recursive depth k and across regimes featuring implicit signaling (like coded communication, steganographic communication)
\end{itemize}

\section{Talk with Jerick}
2 Proposals: 
\begin{enumerate}
    \item Understand what agents can do in open source game theory, using lean:
    \begin{itemize}
        \item Transparency (Master Thesis): First idea (extension of Critch's paper): We want to answer the open questions of the paper. To do so, we want to code in Lean the proof\_checker (or proof\_writer ?) of a specific game (e.g: $\cupod$ vs $\cupod$ in prisoners dilemma) to be able to simulate the expected outcome (in the example: (D, D)). Notion of length of proof (k) important.
        
        Lean might not be the best for this, will have to check. Maybe a llm that does the lean code, or something else?      
        
        \item Partial Transparency: How much info can we give to the agents so that they achieve cooperation? More on simulation with less info.
        \item Pay a cost to simulate: Analyze thea game theory behind the games (this is the math part) and then do some simulations.
    \end{itemize}
    \item Tackle open-source game theory with other tools than Lean:
    \begin{itemize}
    \item Transform the concepts of (hash being hackable $\to$ hash families better) in a gt sense.
    Instead of committing into one strategy, you have a family of strategies for agent A. Then can agent B devise a strategy to not get exploited (partial information setting)?
    \item RL in gt: Instead of just feeding the reward to the agent, we also feed in the weights of the last layer of the opponent → RL training: see if there is better behavior.
    \item Creating a benchmark to see how good agents are at simulating the actions of other agents. What are the different outcomes? Difference between one shot reasoning and recursive reasoning.
\end{itemize}
\end{enumerate}


Material on the Jinesis drive I can read, in folders: presentations, proofs, experiments.
\end{document}